% ═══════════════════════════════════════════════════════════════
% 第九章：结论与政策建议
% ═══════════════════════════════════════════════════════════════

\section{结论与政策建议}

\subsection{主要结论}

本文基于中国A股上市公司约10万公司-季度观测数据，融合财务指标、大语言模型提取的文本语义特征、市场特征、公司治理特征和历史监管信息等多源异构数据，构建了面向监管问询概率预测与可解释归因的综合分析框架。主要发现如下：

第一，传统的财务异常指标具有显著的预测能力，Logistic回归AUC达到0.71。Beneish M-Score、应收增速偏离和经营性现金流质量是三个最重要的财务预测因子。

第二，大语言模型（LLM）提取的文本语义特征提供了显著的增量预测信息。仅文本语义模型即可达到与全量财务指标接近的预测性能（AUC 0.73），且财务指标与文本语义特征的融合在Delong检验中表现出显著的互补效应（$p<0.01$）。这为"软信息"在公司风险预测中的独立价值提供了新的证据。

第三，软信息的增量价值在信息不对称严重的公司中更为凸显。小市值公司（+0.085）和无分析师覆盖公司（+0.091）的增量AUC约为大市值和高覆盖公司的2倍。这一异质性特征支持了信息不对称理论——软信息在公开信息不充分的环境中更有可能发挥信息补充的功能。

第四，LightGBM异质集成全模型的AUC达到0.79，Top-10\%高概率公司覆盖了38.6\%的真实问询样本，表明系统能够有效识别相当一部分在实际中被监管机构问询的重点公司。

第五，基于SHAP归因和LLM生成的五部分预警报告在LLM-as-Judge评估中获得了86.3分的综合质量评分，验证了从概率到证据的三级归因路径的有效性。

\subsection{理论贡献}

本文的理论贡献表现在三个层面。

在信息披露经济学层面，本文首次系统检验了LLM提取的文本语义特征对监管问询的增量预测价值，将Liberti和Petersen（2019）的"软信息—硬信息"分析框架从传统信用风险评估拓展至监管问询预测这一新场景。我们发现软信息的增量价值在信息不对称更严重的公司中更大，这为揭示软信息发挥作用的条件边界提供了定量证据。

在监管经济学层面，本文构建的监管问询概率预测模型本质上是对监管者行为函数的估计和分解。研究结果揭示了监管问询行为的系统选择性——监管者更关注收入质量异常、关联交易风险和信息披露矛盾等特征，这些发现有助于深化对监管问询行为逻辑的理解。

在金融科技方法论层面，本文构建了一个融合传统计量方法、机器学习和生成式人工智能的可解释预测框架，该框架在保持预测性能的同时，通过SHAP归因和LLM报告生成实现了高水平的可解释性，为金融领域AI应用的"可解释性合规"提供了一个可参考的技术方案。

\subsection{政策建议}

基于上述研究结论，本文提出以下政策建议。

\textbf{第一，建立基于多源信息融合的智能化问询筛选系统。}研究发现LLM语义特征的加入能够显著提升问询预测的准确性。建议交易所在日常信息披露审查中，将文本语义分析作为传统财务指标审查的补充手段，构建"财务指标+文本语义+市场信号"三位一体的预警体系，进一步提升问询函的靶向性和监管资源的配置效率。

\textbf{第二，重点关注信息不对称严重的中小市值公司。}实证结果表明，软信息在小市值和低分析师覆盖公司中的增量价值最为突出。这类公司由于市场关注度较低，信息透明度不足，往往是信息披露问题的"高发区"。建议监管机构在资源有限的条件下，优先将文本语义分析工具应用于中小市值公司的风险筛查。

\textbf{第三，将可解释的AI归因纳入监管科技合规体系。}本文构建的三级归因机制实现了从预测概率到原文证据的完整追溯，满足了金融场景对AI决策可解释性、可审计和可复核的要求。建议监管部门在推进金融科技赋能监管时，将"可解释性"作为AI辅助工具的核心考量标准，确保AI输出的每一项风险结论都可以追溯到具体的证据来源。

\textbf{第四，鼓励上市公司提升信息披露质量，尤其是文本披露的充分性和可理解性。}研究发现公告文本的风险要素计数与问询概率显著正相关，意味着信息越不充分、越含糊的公告越可能触发监管问询。上市公司应重视MD\&A等文本披露内容的信息含量，避免形式化、模板化的披露方式。

\subsection{研究局限与展望}

本研究存在以下局限，也为未来研究提供了可拓展的方向。第一，受数据可得性限制，本文使用的大语言模型为通用版本的Qwen3-7B模型，未在金融领域数据进行专门的微调训练。如果模型能够在上市公司公告数据库上进行领域适应（Domain Adaptation），其抽取的特征质量有望进一步提升。第二，本文的相似案例检索主要基于文本向量相似度，尚未充分利用问询函之间的法律逻辑关联和监管意图相似性。引入知识图谱推理和因果推断方法有望提升案例匹配的精确度。第三，本文的预测框架主要面向二分类（"是否被问询"），尚未对问询函的类型和严重程度进行细粒度预测。将问询函按类型和问题的数量、性质进行进一步细分，将是更有业务应用价值的研究方向。第四，也是最重要的，本文的实证结果基于虚拟数据集（mock data）产生，后续需要替换为真实数据集后方可形成最终结论。